\chapter{Optimization Algorithms}

\vspace*{1cm}

Optimization algorithms play a central role in the training of neural networks.
During fine-tuning, the model parameters are updated in order to minimize the loss
function. The choice of optimizer can influence convergence speed, training
stability, memory consumption, and final model performance. This chapter presents
the optimization methods studied in this project.

\section{Stochastic Gradient Descent}

Stochastic Gradient Descent is one of the most fundamental optimization
algorithms in machine learning. It updates the model parameters using the gradient
of the loss function computed on mini-batches of training data.

\vspace*{0.5cm}

The update rule can be written as:

\begin{equation}
    \theta_{t+1} = \theta_t - \eta \nabla_{\theta} L(\theta_t)
\end{equation}

where \(\theta_t\) represents the model parameters at step \(t\), \(\eta\) is the
learning rate, and \(\nabla_{\theta} L(\theta_t)\) is the gradient of the loss
function with respect to the parameters.

\vspace*{0.5cm}

SGD is simple and memory-efficient, but it can converge slowly and may be
sensitive to the choice of learning rate. In transformer fine-tuning, plain SGD
is generally less stable than adaptive methods, but it remains useful as a
baseline optimizer.

\section{SGD with Momentum}

SGD with momentum improves classical SGD by accumulating part of the previous
update direction. This helps accelerate learning in consistent gradient
directions and reduces oscillations during optimization.

\vspace*{0.5cm}

The update rule is:

\begin{align}
    v_{t+1} &= \mu v_t + \nabla_{\theta} L(\theta_t) \\
    \theta_{t+1} &= \theta_t - \eta v_{t+1}
\end{align}

where \(v_t\) is the velocity term and \(\mu\) is the momentum coefficient.

\vspace*{0.5cm}

Momentum can improve convergence compared with plain SGD. However, it still uses
a global learning rate and does not adapt the learning rate independently for
each parameter. This can make it less flexible for large transformer models.

\section{Adam Optimizer}

Adam, which stands for Adaptive Moment Estimation, is one of the most widely used
optimizers in deep learning. It combines ideas from momentum and adaptive learning
rates. Adam keeps an exponential moving average of both the gradients and the
squared gradients.

\vspace*{0.5cm}

The main update equations are:

\begin{align}
    m_t &= \beta_1 m_{t-1} + (1 - \beta_1) g_t \\
    v_t &= \beta_2 v_{t-1} + (1 - \beta_2) g_t^2
\end{align}

where \(g_t\) is the gradient at step \(t\), \(m_t\) is the first moment estimate,
and \(v_t\) is the second moment estimate.

\vspace*{0.5cm}

Adam is generally well suited for transformer fine-tuning because it adapts the
learning rate for each parameter and often converges faster than SGD-based
methods. For this reason, it is commonly used as a strong baseline in language
model training.

\section{L-BFGS}

L-BFGS, or Limited-memory Broyden-Fletcher-Goldfarb-Shanno, is a quasi-Newton
optimization algorithm. Unlike SGD and Adam, which mainly rely on first-order
gradients, L-BFGS approximates second-order curvature information in order to
improve the search direction.

\vspace*{0.5cm}

The method does not store the full Hessian matrix. Instead, it keeps a limited
history of previous parameter and gradient differences. This makes it more memory
efficient than full Newton methods while still using curvature information.

\vspace*{0.5cm}

In this project, L-BFGS is included because it is important from an optimization
course perspective. However, it is usually expensive for transformer fine-tuning.
It requires closure-based optimization steps and may perform several evaluations
of the loss during a single update. This can increase training time compared with
optimizers such as Adam.

\section{Newton-CG}

Newton-CG, or Newton Conjugate Gradient, is a second-order optimization method.
It is based on Newton's method and uses conjugate gradient techniques to solve the
linear system involving the Hessian matrix.

\vspace*{0.5cm}

The ideal Newton update can be expressed as:

\begin{equation}
    \theta_{t+1} = \theta_t - H^{-1} \nabla_{\theta} L(\theta_t)
\end{equation}

where \(H\) is the Hessian matrix of second derivatives.

\vspace*{0.5cm}

Although Newton-CG can be powerful for some optimization problems, it is not
practical for standard fine-tuning of transformer models in this project. The
method requires Hessian-vector products and second-order computations that are
expensive for models with many parameters. For this reason, Newton-CG is kept for
theoretical discussion only and is not implemented in the training pipeline.

\section{Comparison Perspective}

The selected optimizers represent different families of optimization methods.
SGD is a simple first-order baseline. SGD with momentum improves the direction of
updates by using past gradients. Adam introduces adaptive learning rates and is
commonly used in transformer training. L-BFGS brings a quasi-Newton perspective,
while Newton-CG represents a theoretical second-order method.

\vspace*{0.5cm}

The comparison between these optimizers is important for understanding the
tradeoff between training quality and computational efficiency. In the context of
Edge AI, the best optimizer is not only the one that gives the lowest validation
loss, but also the one that offers a good balance between performance, training
time, memory usage, and inference efficiency.
